% Figure: the tanh-sinh (double-exponential) change of variables maps
% [-1,1] to the whole real line so that the transformed integrand decays
% doubly exponentially. We plot the node density: x = tanh((pi/2) sinh t)
% as a function of t, showing how uniform t-spacing clusters nodes near
% the endpoints +-1 where an endpoint singularity lives.
\begin{figure}[htbp]
\centering
\begin{tikzpicture}
\begin{axis}[
  width=11cm, height=6cm,
  xlabel={$t$ (uniform grid)}, ylabel={$x = \tanh(\frac{\pi}{2}\sinh t)$},
  xmin=-3.2, xmax=3.2, ymin=-1.08, ymax=1.08,
  domain=-3.2:3.2, samples=300,
  axis lines=middle, grid=both,
  major grid style={engblack!12}, font=\small,
  legend pos=south east, legend style={font=\footnotesize},
]
\addplot[engblue, very thick] {tanh((pi/2)*sinh(x))};
\addlegendentry{$x(t)$ map}
% uniform-in-t nodes mapped to x, showing endpoint clustering
\addplot[only marks, mark=*, engred, mark size=1.3pt]
  coordinates {(-3,-1)(-2.5,-1)(-2,-0.99996)(-1.5,-0.9967)(-1,-0.9476)
  (-0.5,-0.7026)(0,0)(0.5,0.7026)(1,0.9476)(1.5,0.9967)(2,0.99996)
  (2.5,1)(3,1)};
\addlegendentry{equispaced $t$ nodes}
\end{axis}
\end{tikzpicture}
\caption[The tanh-sinh node map]{The tanh-sinh (double-exponential)
change of variables $x = \tanh(\frac{\pi}{2}\sinh t)$ maps a uniform
grid in $t$ to nodes that cluster doubly exponentially toward the
endpoints $x = \pm 1$ (red dots). The transformed integrand carries a
weight $\frac{\mathrm{d}x}{\mathrm{d}t}$ that decays like
$\exp(-\frac{\pi}{2}\cosh t)$, so an endpoint singularity in the
original integrand is crushed by a weight that vanishes faster than the
singularity grows. Trapezoidal integration in $t$ over a modest range,
truncated where the weight underflows, then attains near-machine
accuracy on integrands with integrable endpoint singularities that
defeat every fixed-order polynomial rule. The map is the workhorse of
arbitrary-precision quadrature.}
\label{fig:vol02:ch06:tanh-sinh}
\end{figure}
